% Problem-section file following ../_template.tex.
\section{Unconditional classical verification with one quantum prover}
\label{sec:problem-39}

\paragraph{Problem.}
Does every language $L\in\mathsf{BQP}$ admit a single-prover interactive proof
with a fully classical verifier, an efficient quantum honest prover, and
information-theoretic soundness?  Precisely, require a probabilistic classical
polynomial-time verifier exchanging only classical messages with one prover;
a uniform quantum polynomial-time honest prover; acceptance probability at
least $2/3$ for every yes-instance when the honest prover is used; and
acceptance probability at most $1/3$ for every no-instance against every
prover, including a computationally unbounded one.  Equivalently, can one
classically verify an arbitrary efficient quantum computation with one
efficient quantum prover, no quantum capability for the verifier, and no
cryptographic assumption?

\subsection*{Status}

Unsolved

\subsection*{Source}

Aharonov, Ben-Or, Eban, and Mahadev explicitly ask whether the verifier's
remaining quantum register can be eliminated from an unconditional
single-prover verification protocol
\sourcecite{ref:p39-aharonov}{ABE+17}.

\subsection*{Progress}

\begin{itemize}
  \item The inclusions
  $\mathsf{BQP}\subseteq\mathsf{PSPACE}=\mathsf{IP}$ give classical-verifier
  interactive proofs for $\mathsf{BQP}$, but do not ensure that the honest
  prover is implementable in quantum polynomial time
  \sourcecite{ref:p39-adleman}{ADH97},
  \sourcecite{ref:p39-shamir}{Sha92}.

  \item Unconditional protocols with an honest $\mathsf{BQP}$ prover are
  known when the verifier retains a constant-size quantum register or
  exchanges a small number of qubits
  \sourcecite{ref:p39-aharonov}{ABE+17}.  They do not meet the fully classical
  verifier and classical-message requirements.

  \item A fully classical single-verifier protocol is known under the
  learning-with-errors assumption
  \sourcecite{ref:p39-mahadev}{Mah18}.  Its soundness is computational against
  efficient quantum provers, rather than information theoretic against
  arbitrary provers.

  \item Fully classical and unconditional verification is possible with
  multiple noncommunicating entangled provers
  \sourcecite{ref:p39-reichardt}{RUV13}.  This changes the single-prover model
  required in the problem statement.
\end{itemize}

\subsection*{References}

\begin{enumerate}[label={},leftmargin=4.8em,labelsep=0.5em]
  \footnotesize
  \item[\textup{[ADH97]}]\label{ref:p39-adleman}
  L. M. Adleman, J. DeMarrais, and M.-D. A. Huang,
  ``Quantum Computability,'' \emph{SIAM Journal on Computing} \textbf{26},
  1524--1540 (1997).
  \href{https://doi.org/10.1137/S0097539795293639}%
       {doi:10.1137/S0097539795293639}.

  \item[\textup{[Sha92]}]\label{ref:p39-shamir}
  A. Shamir, ``IP $=$ PSPACE,'' \emph{Journal of the ACM} \textbf{39},
  869--877 (1992).
  \href{https://doi.org/10.1145/146585.146609}%
       {doi:10.1145/146585.146609}.

  \item[\textup{[ABE+17]}]\label{ref:p39-aharonov}
  D. Aharonov, M. Ben-Or, E. Eban, and U. Mahadev,
  ``Interactive Proofs for Quantum Computations,'' arXiv:1704.04487 (2017).
  \href{https://arxiv.org/abs/1704.04487}{arXiv:1704.04487}.

  \item[\textup{[Mah18]}]\label{ref:p39-mahadev}
  U. Mahadev, ``Classical Verification of Quantum Computations,'' in
  \emph{Proceedings of the 59th IEEE Symposium on Foundations of Computer
  Science}, 259--267 (2018).
  \href{https://doi.org/10.1109/FOCS.2018.00033}%
       {doi:10.1109/FOCS.2018.00033};
  \href{https://arxiv.org/abs/1804.01082}{arXiv:1804.01082}.

  \item[\textup{[RUV13]}]\label{ref:p39-reichardt}
  B. W. Reichardt, F. Unger, and U. Vazirani,
  ``Classical Command of Quantum Systems,'' \emph{Nature} \textbf{496},
  456--460 (2013).
  \href{https://doi.org/10.1038/nature12035}{doi:10.1038/nature12035};
  \href{https://arxiv.org/abs/1209.0449}{arXiv:1209.0449}.
\end{enumerate}

\subsection*{Comment}

Section~1.8 of \sourcecite{ref:p39-aharonov}{ABE+17} explicitly asks whether
the verifier's remaining quantum register can be eliminated.  The unresolved
conjunction is one prover, a fully classical verifier, a $\mathsf{BQP}$ honest
prover, and unconditional soundness; relaxing any one of these requirements
leads to known protocols.

\subsection*{Tag}

Interactive proof systems; Delegated quantum computation;
Verification and validation; Quantum complexity theory

\subsection*{ID}

\texttt{op\_3770ad932d54692f}
